\documentclass[11pt]{article}
\usepackage[utf8]{inputenc}
\usepackage{amsmath,amssymb}
\usepackage{authblk}
\usepackage[margin=1in]{geometry}
\usepackage[hidelinks]{hyperref}
\usepackage{xurl}
\usepackage{setspace}

\title{Deterministic RealQM Alpha Decay:\\
Coulomb-Barrier Tunnelling and the Geiger--Nuttall Law}
\author[1]{Claes Johnson}
\affil[1]{KTH Royal Institute of Technology, SE-100 44 Stockholm, Sweden \\
\texttt{claesjohnson@gmail.com}}
\date{\today}

\begin{document}
\maketitle

\begin{abstract}
Alpha decay is the clean case for a charge-density account of radioactivity. It is a genuine \emph{two-body}
decay --- daughter nucleus plus an alpha particle of charge $+2$ --- driven by tunnelling through the
daughter's \emph{Coulomb} barrier, with no weak interaction, \emph{no neutrino}, and a \emph{discrete,
monoenergetic} alpha line whose sharpness is itself the proof that no third body is emitted. Everything the
process needs --- extended charge, a Coulomb barrier, two-body kinematics --- lies inside RealQM; nothing of
the chargeless carrier that a companion analysis of beta decay had to concede is present here. We make two
points. First, quantitatively: a single WKB Coulomb-barrier (Gamow) calculation reproduces the
Geiger--Nuttall law --- $\log_{10}t_{1/2}$ linear in $Q^{-1/2}$ --- across the twenty-five orders of
magnitude that alpha lifetimes span, matching measured half-lives from $^{232}$Th ($1.4\times10^{10}$~yr) to
$^{212}$Po ($0.3\,\mu$s) within a factor of a few, with the Geiger--Nuttall slope $158$ against the measured
$153$. The timing of alpha decay is barrier tunnelling of an extended charge, entirely a charged-sector
phenomenon. Second, foundationally: in time-dependent RealQM the tunnelling escape is \emph{deterministic},
triggered by a phase coincidence among the charge domains rather than by irreducible chance, and the
exponential, memoryless decay law emerges as the ensemble statistics of that coincidence over unknown
initial phases. Alpha decay is thus a deterministic Coulomb process throughout, with the neutrino nowhere in
sight.
\end{abstract}

\setstretch{1.05}

\section{Why alpha decay is the clean case}
\label{sec:clean}

Radioactive decay puts two questions to a theory: \emph{when} a nucleus decays, and \emph{what carries} the
released energy and momentum. For beta decay the second question forces a chargeless carrier --- the
antineutrino, which takes on average some $60\%$ of the released energy and the momentum imbalance, and
which lies outside a charge-density ontology~\cite{BetaCompanion}. Alpha decay is different in exactly the
way that matters.

Alpha decay is $\mathrm{parent}\to\mathrm{daughter}+\alpha$, the alpha being a $^4$He nucleus of charge
$+2$. It is a genuine \emph{two-body} decay: no leptons are created, no weak interaction acts, and no
neutrino is emitted. The experimental signature is decisive: alpha spectra are \emph{discrete}, each branch
emitting the alpha at a sharp, fixed energy (the alpha carrying $\approx(A{-}4)/A$ of the release, the heavy
daughter the small recoil). A two-body decay from rest is forced by conservation to send the products
back-to-back with a \emph{fixed} energy split; a continuous alpha spectrum would signal a hidden third body,
as the continuous \emph{electron} spectrum does in beta decay. The alpha line is sharp. That sharpness is
the direct evidence that alpha decay has \emph{no} missing carrier.

Every ingredient is then inside RealQM. The alpha is a charged, extended density; the barrier it must cross
is the \emph{Coulomb} repulsion of the daughter; the kinematics is two-body. A RealQM two-body decay comes
out exactly back-to-back and monoenergetic~\cite{BetaCompanion}, which is \emph{wrong} for beta decay but
\emph{right} for alpha. Where beta decay exposed the boundary of the charge-density picture, alpha decay
sits squarely within it.

\section{The timing: Coulomb-barrier tunnelling reproduces Geiger--Nuttall}
\label{sec:gamow}

The alpha is held inside the nucleus by the short-range attraction and confined by the Coulomb barrier
$V(r)=2Z_d e^2/r$ of the daughter (charge $Z_d$) outside the nuclear radius $R$. Classically trapped, it
escapes only by tunnelling --- the picture of Gamow and of Condon and Gurney~\cite{Gamow1928,Condon1928}.
The WKB barrier factor for an alpha of energy $Q$ is
\begin{equation}
2G=\frac{2}{\hbar}\sqrt{2m_\alpha}\int_R^{b}\!\sqrt{V(r)-Q}\;dr
=\frac{2}{\hbar}\sqrt{2m_\alpha}\,\sqrt{Q}\,b\,\Big[\arccos\sqrt{x}-\sqrt{x(1-x)}\Big],
\label{eq:gamow}
\end{equation}
with outer turning point $b=2Z_d e^2/Q$ and $x=R/b$, and the decay rate is $\lambda=f\,e^{-2G}$ for an
assault frequency $f$. Because $2G$ scales as $Q^{-1/2}$, this yields the Geiger--Nuttall law
\begin{equation}
\log_{10}t_{1/2}=\frac{a}{\sqrt{Q}}+b .
\label{eq:gn}
\end{equation}

Evaluating \eqref{eq:gamow} for a series of even--even emitters, with nuclear radius
$R=1.20\,(A_d^{1/3}+4^{1/3})$~fm and no fitted parameters, reproduces measured half-lives across the whole
Geiger--Nuttall range (Table~\ref{tab:gn}): from $^{232}$Th at $Q=4.08$~MeV and $t_{1/2}\approx1.4\times
10^{10}$~yr to $^{212}$Po at $Q=8.95$~MeV and $t_{1/2}\approx0.3\,\mu$s --- \emph{twenty-five orders of
magnitude} --- each within a factor of a few. The fitted Geiger--Nuttall slope is $a=158$ for the model
against $a=153$ for the data, and the model spans $25.1$ decades against the measured $24.2$.

\begin{table}[h]
\centering
\begin{tabular}{lrrr}
\hline
nuclide & $Q$ (MeV) & $t_{1/2}$ exp. (s) & $t_{1/2}$ Gamow (s) \\
\hline
$^{232}$Th & 4.083 & $4.4\times10^{17}$ & $9.9\times10^{17}$ \\
$^{238}$U  & 4.270 & $1.4\times10^{17}$ & $3.1\times10^{17}$ \\
$^{226}$Ra & 4.871 & $5.1\times10^{10}$ & $6.4\times10^{10}$ \\
$^{238}$Pu & 5.593 & $2.8\times10^{9}$  & $2.1\times10^{9}$  \\
$^{220}$Rn & 6.405 & $5.6\times10^{1}$  & $5.5\times10^{1}$  \\
$^{212}$Po & 8.954 & $3.0\times10^{-7}$ & $8.7\times10^{-8}$ \\
\hline
\end{tabular}
\caption{Alpha half-lives from the parameter-free WKB Coulomb-barrier calculation \eqref{eq:gamow} against
experiment, spanning $\sim25$ orders of magnitude.}
\label{tab:gn}
\end{table}

The lesson is that the timing of alpha decay is fixed by \emph{the Coulomb barrier the extended alpha charge
must cross} --- a purely charged-sector, electrostatic quantity. This is not a new dynamical law; it is the
Gamow picture. What matters for the present programme is that the entire rate, over twenty-five decades,
follows from the Coulomb barrier alone, with neither a strong nor a weak coupling constant nor any carrier
beyond the alpha and the daughter.

\section{The escape is deterministic}
\label{sec:det}

Standard quantum mechanics reads the individual alpha's escape as \emph{irreducibly random}: the barrier
sets a probability per unit time, but \emph{when} a given nucleus decays is held to be uncaused. Time-%
dependent RealQM offers a different reading, the same one that governs beta-decay timing~\cite{BetaCompanion}.
Each charge domain of the configuration carries a phase $\theta_k(t)=-E_kt/\hbar$ fixed by its computed
energy $E_k$, and the alpha crosses the barrier at the first instant when its phase relative to the other
domains lies jointly within an escape window --- a \emph{geometric coincidence} of phases. Given the initial
phases the escape time is determined; the only randomness is ignorance of those phases. Averaged over them,
the coincidences of incommensurate phase-beats form a stationary point process, so the survival probability
decays at a constant hazard rate,
\begin{equation}
S(t)\simeq e^{-\lambda t},\qquad
\lambda\sim\Big(\tfrac{\Delta}{2\pi}\Big)^{M}\langle|\omega_m|\rangle,\quad \omega_m=(E_\alpha-E_m)/\hbar ,
\label{eq:S}
\end{equation}
an exponential, \emph{memoryless} law with a definite half-life. The memorylessness usually read as the
fingerprint of fundamental chance here \emph{emerges} from deterministic phase dynamics as the ensemble
statistics of a coincidence.

That the phase gating is real, and not an imposed rule, is shown by a direct solve. A particle held in a
metastable well behind a barrier --- the Gamow configuration itself --- prepared in a superposition of two
quasi-bound levels $E_1,E_2$ and evolved by a unitary integration of the time-dependent equation with an
absorbing outer layer, escapes with a transmitted flux that is smooth for a single level (ordinary
tunnelling) but \emph{oscillates}, in gated bursts, at the beat frequency $(E_2-E_1)/\hbar$ for the
two-level superposition, to within a few percent, with no escape rule imposed~\cite{PhaseSim}. The relative
phase of the trapped components opens and closes the escape channel. Applied to the alpha behind its Coulomb
barrier, this is a deterministic account of \emph{when} the alpha leaves, with the Gamow factor
\eqref{eq:gamow} setting the scale of the rate and the phase coincidence setting its trigger.

\section{Conclusion}
\label{sec:concl}

Alpha decay is a deterministic Coulomb process from beginning to end. It is two-body, so its products emerge
back-to-back with a discrete, monoenergetic alpha --- the kinematics a RealQM two-body solve produces, and
the direct evidence that no neutrino is emitted. Its timing is barrier tunnelling of the extended alpha
charge: one parameter-free Coulomb-barrier calculation reproduces the Geiger--Nuttall law across
twenty-five orders of magnitude in half-life. And its randomness is only apparent: in time-dependent RealQM
the escape is triggered by a phase coincidence among charge domains, deterministic given the initial phases,
with the exponential law emerging as ensemble statistics. Where beta decay marks the boundary of the
charge-density picture --- the chargeless neutrino carrying energy and momentum the theory cannot supply
\cite{BetaCompanion} --- alpha decay lies wholly within it: charged throughout, Coulomb throughout,
deterministic throughout, with no neutrino anywhere.

\begin{thebibliography}{9}
\bibitem{Gamow1928} G.~Gamow, \emph{Zur Quantentheorie des Atomkernes}, Z. Phys. \textbf{51}, 204 (1928).
\bibitem{Condon1928} R.~W.~Gurney and E.~U.~Condon, \emph{Wave Mechanics and Radioactive Disintegration},
Nature \textbf{122}, 439 (1928).
\bibitem{RealNucleus} C.~Johnson, \emph{RealNucleus: Nuclear Binding as Dual Coulomb Confinement, without a
Strong or Weak Force}, manuscript, 2026.
\bibitem{RealQM} C.~Johnson, \emph{Quantum Mechanics as Multiphase 3D Continuum Mechanics}, manuscript,
2026; \url{https://claes542.github.io/RealMolecule/gallery.html}.
\bibitem{ComplexRealQM} C.~Johnson, \emph{Complex Time-Dependent RealQM: Currents, Radiation, and Magnetism
from the Real-Space Continuum}, manuscript, 2026.
\bibitem{BetaCompanion} C.~Johnson, \emph{Deterministic RealQM Beta Decay: A Phase-Coincidence Half-Life},
manuscript, 2026 (companion paper).
\bibitem{PhaseSim} C.~Johnson, \emph{Phase-gated escape of a metastable state (interactive/numerical
demonstration)}, 2026; \url{https://claes542.github.io/RealMolecule/gallery.html}.
\end{thebibliography}

\end{document}
